\documentclass[10pt,a4paper]{amsart}
\usepackage[margin=19mm]{geometry}
\usepackage{amsmath,amssymb,mathtools}
\usepackage[scr=rsfs,cal=euler]{mathalfa}
\usepackage{xfrac}
\usepackage[unicode,hidelinks,bookmarks=false]{hyperref}
\newcommand{\ie}{i.e.\ }
\newcommand{\cf}{cf.\ }
\newcommand{\resp}{respectively\ }
\newcommand{\Subsection}{\subsection}
\providecommand{\nobreakdash}{\nobreak\hbox{-}}
\setlength{\emergencystretch}{2em}
\multlinegap0pt
\usepackage{amssymb}
\usepackage{xcolor}
\usepackage{tikz}
\usepackage{bm}
\usepackage{mathtools}
\usepackage[shortlabels]{enumitem}
\usepackage{csquotes}
\usepackage{float}
\newtheorem{proposition}{Proposition}
\newcommand{\Ht}{\widetilde{H}}
\newcommand{\ZZ}{\mathbb{Z}}
\newcommand{\cF}{\mathcal{F}}
\newcommand{\cK}{\mathcal{K}}
\newcommand{\cR}{\mathcal{R}}
\newcommand{\del}{\partial}
\title{Constructions of symplectic forms on 4-manifolds}
\author{Peter Lambert-Cole}
\date{Source: arXiv:2403.05512v2 (CC BY 4.0)}
\begin{document}
\maketitle

\noindent\emph{Source and licence.} Constructions of symplectic forms on 4-manifolds. Version arXiv:2403.05512v2 (CC BY 4.0), distributed by arXiv under the Creative Commons Attribution 4.0 International licence, http://creativecommons.org/licenses/by/4.0/, as recorded in the arXiv licence metadata for this exact version and shown on the abstract page https://arxiv.org/abs/2403.05512v2. Full licence notice: \texttt{licenses/arxiv-2403.05512-CC-BY-4.0.txt}.\par\smallskip

\noindent\textit{Editorial scope.} Counted unit: the article's proposition prop:all-line-bundles (source0.tex lines 792--795). The article's complete original proof of it is delivered with it (source0.tex lines 797--827, one \texttt{proof} environment of 31 lines). No mathematics is written, repaired or re-derived: the statement and the proof are the article's own lines, in the article's own notation, and no retained line is substituted or re-flowed.  No further source-local context is needed: every cross-reference of every retained line resolves inside the two delivered spans.  Nothing was omitted from any retained span, so the package carries no omission record.  No retained line contains a citation, so the wrapper carries no bibliography and no bibliography entry.  No retained line contains a figure environment, an image-inclusion command or drawing code, so no image is delivered and the package declares no asset dependency.  Editorial formatting only, in addition: an AMS \texttt{amsart} preamble that restates the article's own declarations and macros used by the retained lines, namely the environments \texttt{proposition} on separate counters, together with the standard packages those lines need.  The retained environments print in this document as Proposition 1 in order of appearance, whereas the article prints the same items with the numbering of its own full document.  The complete native source of the article as distributed in the arXiv e-print for this exact version is bundled unchanged as \texttt{sources/155801/source0.tex}.
\begin{proposition}
\label{prop:all-line-bundles}
    Let $(X,\omega)$ be a symplectic 4-manifold with rational symplectic form. The branch locus $\cR$ in the construction can be chosen such that every homology class $A \in H_2(X;\ZZ)$ with positive symplectic area is represented by a surface that is geometrically transverse to the trisection of $X$.
\end{proposition}

\begin{proof}
There exists a surjective map
\[\pi_1(\widetilde{H}_1,p_0) \rightarrow H_1(\widetilde{H}_1) \rightarrow H_2(X)\]
The first is given by the abelianization, while the second is due to the fact that $H_1(\widetilde{H}_1) \cong H_2(\widetilde{H}_1,\del \widetilde{H}_1)$ and there is a handle decomposition of $X$ such that compressing disks in $\widetilde{H}_1$ are precisely the cores of the 2-handles.

Let $A \in H_2(X;\ZZ)$ be a class of positive symplectic area and let $\zeta$ be a loop based at $p$ such that $[\zeta] \in \pi_1(\widetilde{H}_1,p)$ is a lift of $A$.  By assumption, $\zeta$ is nullhomologous in $Y_1,Y_3$ and therefore bounds slice surfaces in $Z_1,Z_3$ whose union is a surface representing $A$.  

As shown in the following subsection, there is a decomposition
\[ \zeta = l_1^n \mu_1 \rho_1 \cdots \mu_k \rho_k\]
where
\begin{enumerate}
    \item $l_1$ is loop in $\Ht_1$ transverse to the foliation determined by $\beta_1$
    \item each loop $\mu_i$ is tangent to the leaf of $\cF_1$ through the basepoint $p$,
    \item each $\rho_i$ is represented by a loop on $\Sigma$ that is homotopic in $\Ht_2$ to a loop in the leaf of $\cF_2$ through $p$ and that is contractible in $\Ht_3$.
\end{enumerate}

In a neighborhood of $p$, the defining 1-forms $\beta_1,\beta_2,\beta_3$ exact with local primitives $g_1,g_2,$ and $g_3 = -g_1 - g_2$.  Pick $2k$ points $p_1,\dots,p_{2k}$ is a small neighborhood of the basepoint $p$ such that
\begin{align*}
    g_1(p_i) > g_1(p_{i-1}) & &&\text{ for $i$ odd} \\
    g_2(p_i) > g_2(p_{i-1}) & && \text{ for $i$ even}\\
    g_3(p_i) > g_3(p_{i-1}) & &&\text{ for $i$ even}
\end{align*}

Note that since the points $\{p_i\}$ are in an arbitrary neighborhood of $p$, there is a unique identification between homotopy classes of loops based at $p$ and homotopy classes of paths from $p_j$ to $p_{j'}$ for any $j,j'$.

To build the surface, connect $p_{2i-1}$ to $p_{2i}$ by arcs of $\tau_2$ and $\tau_3$ homotopic to $\rho_i$ and connect $p_{2i}$ to $p_{2i+1}$ be an arc homotopic to $\mu_{i+1}$.  Finally, connect $p_{2k}$ to $p_1$ by a loop homotopic to $l_1^n \mu_1$.  If $n > 0$, then each arc of $\tau_{\lambda}$ can be assumed geometrically transverse to the foliation $\cF_{\lambda}$.  

From Stokes's Theorem, it follows that
\[\int_{\cK} \omega = \sum_{\lambda = 1}^3 \int_{\tau_{\lambda}} \beta_{\lambda}\]
Since all of the component arcs of $\tau_2,\tau_3$ are perturbations of flat paths, their contribution is $O(\epsilon)$, whereas the contribution of $\tau_1$ to his is approximately $n \int_{l_1} \beta_1$.  Therefore, the surface has positive symplectic area if and only if $n > 0$.
\end{proof}

\end{document}
